\section{Empirical and theoretical implications}
\subsection{Choice-homophily of network contacts based on pre-existing cultural dispositions}
Standard network theories of taste transmission draw on the notion of ``homophily'' (i.e. the empirical generalization that close friends tend to be similar in sociodemographic characteristics) to explain the fact that cultural tastes and practices tend to cluster in certain areas of social space \citep{mark03}.  From the standard point of view homogeneous social connections based on socio-demographic similarity pre-exist individual patterns of taste and this explains why cultural practices ``lump'' in certain areas of social space.  That is, the logic of network-based theories leads to them to attempt to explain the fact that cultural practices occupy delimited niches in social space by pointing to two interrelated processes:  (1) tastes tend to be associated with social background and (2) individuals tend to associate with those who are socio-demographically similar.  

Yet, there is sufficient qualitative \citep{cardon_granjon05, long03, wali_etal02} and quantitative \citep{lizardo06a, vaisey_lizardo10, wimmer_lewis10} evidence to suggest that just in the very same that tastes and subcultural preferences and values are transmitted through pre-existing network ties, new social ties are formed on the basis of pre-existing cultural dispositions \citep{lazarsfeld_merton54}.  This is the obverse of the traditional network-theoretic emphasis on the formation of new cultural tastes on the basis of preexisting social ties.  In fact, given the relative stability of cultural tastes in comparison to network connections, it is likely that this last process may in fact occur more often that the former process \citet{lizardo06a}.  In this respect (and in contradiction to standard network accounts), the ``lumpiness'' of cultural tastes in social space may be actually instead be the \emph{result} of an active---but not necessarily conscious---process of contact selection and relationship fine-tuning keyed around cultural compatibility \citep{lazarsfeld_merton54, vaisey_lizardo10, dimaggio93b}.

 This proposition is concordant with Burt's  observation of the ``liability of newness'' of recently formed ties, and with recent studies of changes in social networks of young adolescents and young adults \citep{bidart_degenne05, bidart_lavenu05, cardon_granjon05}.  If an initial ``feeling-out'' process (keyed to cultural compatibility) is characteristic of newly formed relationships, and if culturally incompatible contacts are likely to drop the relationship at early signs of incompatibility, then the higher mortality hazard of new relationships is readily explicable \citep{vaisey_lizardo10}.  Notice that if this is correct, then sociodemographic homophily could in the majority of cases be a by-product of cultural compatibility and not the other way around (this cultural compatibility may of course be itself due to socio-demographic of the family of origin [Bourdieu 1984; Illouz 1998]), as is strongly suggested in recent study of the cultural-underpinnings of racial homophily using dynamic network data \citep{wimmer_lewis10} and of the changing composition of adolescents networks using dynamic panel data \citep{vaisey_lizardo10}.  This is an insight that was always implicit in the theoretical logic of dynamic simulations of network formation based on shared cultural knowledge \citep{carley91, dimaggio93b}, but which has been occluded in recent theory and research due to the one sided emphasis put on the socio-structural determination of shared tastes.  
 
\subsection{Class-based heterogeneity in contact choice and structural determination of tastes}
Another important implication of the above considerations is the fact that the extent to which contacts are subject to a process of active selection based on taste itself varies according to the individual's command of cultural, educational and socio-economic resources.  Put simply, high cultural capital individuals, in particular those employed in prestigious occupations that demand of extensive degrees of cognitive and cultural skill not only have larger and sparser social networks \citep{lizardo06a}, but the proportion of their networks that is composed of constantly changing, loosely bounded social ties is larger and more heterogeneous than that of culturally disadvantaged individuals employed in routine, low-complexity occupations \citep{coser91, fischer82}.   In this respect, the endogenous effect of cultural capital in both the ability to select among different social ties and the ability to reach across the social structure to form those ties, is itself partially determined by position in social space. 
 
This is important, because the dominant way that ``social space'' has been conceptualized by network and organizational theorists \citep{hannan_etal03, mcpherson04} in relation to culture has been---following \citet{blau77a}---as a multidimensional manifold composed of variables such as age and occupational prestige.  The sociodemographic niche for a particular cultural or social form can then be given a metric definition in terms of these variables.  Social distance among two points in these space is conceived as being proportional to sociometric distance in social networks, that as being a proxy for the probability that two individuals will interact in real life \citep{mark98a}.  This presumes that the socio-demographic covariates that define social space are strictly exogenous (i.e. they do not interact in the statistical sense) with the propensity to deploy cultural knowledge reaped from consuming the cultural form in question in interaction to form new social relationships with persons located at relatively distant locations in social space.   Traditional network models also assume that there is no interaction between privilege locations in social space and the relative propensity to acquire or learn about specific cultural forms from proximate social contacts.  This process is seen as being a generic dynamic applicable to all individuals regardless of their socio-demographic profile.  

From the point of view of the culture-conversion model on the other hand the axes of social space as defined above are clearly composed of variables some of which (i.e. educational attainment) are themselves implicated in the extent to which tastes may be driven by social network dynamics.  Thus, from this point of view, occupying a privileged position in social space plays an important role in the degree to which individuals are able to loosen the limiting effects of network-based social constraint. This happens due the greater ability of those who posses higher levels of educational and socio-economic resources to deploy a more heterogeneous repertoire of cultural knowledge in everyday interaction.   This allows them to form and sustain social relationships that cut-across social dimensions of stratification and differentiation \citep{dimaggio87}, as well permitting them to create and maintain intimate ties with those who share their knowledge and tastes.  For this reason, the dimensions of ``social space'' cannot be considered strictly exogenous to the cultural contents whose sociodemographic niche they define.  Instead differential command of these cultural contents on the part of individuals may themselves ``warp'' these dimensions, such that individuals located at some privileged locations in this manifold find it easier to traverse a large distance in that space and thus connect to a wide range of others \citep{erickson96, lin_dumin86}. Individuals located in less privileged areas of the space, on the other hand, are socially and culturally distant from most individuals and cannot connect very well across their locally bounded social circles \citep{bryson97}. 

There is plenty evidence that is consistent with these considerations.  For instance \citet{wellman79, wellman_etal88, wellman_wortley90} and \citet{fischer82} show that as socio-economic and educational status increases, individuals are much more likely to have role-segmented, specialized relationships in which a clear differentiation is made between those ties designed for sociability and companionship (usually labeled as ``friends'') and social ties that are the result of non-negotiable statuses such as kinship ties.  These lines of research on personal community networks also show that individuals are much more likely to prefer to obtain companionship from voluntary rather than involuntary ties.  These are the social contacts that the culture conversion model predicts are likely to be formed and sustained through pre-existing forms of cultural compatibility.  These are also the social relationships that are likely to be sustained through interaction rituals in which the symbolic goods coming from artworlds and cultural industries figure as the primary subject of conversation \citep{cardon_granjon05, dimaggio87}.  Finally, these are also the same social ties that are continually undergoing high rates of turnover and being constantly recycled and reformed, and which thus require the recurrent deployment of cultural capital in interaction.  This goes a long way toward explaining the almost insatiable demand for novel symbolic goods among high cultural capital class fractions \citep{douglas_isherwood96}

In essence, for a relatively privileged cultural and economic elites social relationships have come to be partially disembedded from place and kin constraints \citep{dimaggio_mohr85, holt98}, and have taken the form of increasingly longitudinally unstable and thus constantly changing personal community networks\citep{wellman_etal88}.  This increases the chances that differential command of the ability to consume a wide variety of cultural forms becomes a driver of the recurrent formation of new network relations.  This is less likely to be the case for those  individuals who are restricted to a more homogeneous, kin-based, and locally constrained set of core alters who change relatively little over time.  Consistent with this claim, \citet{lizardo06a} shows that after holding cultural capital constant, the positive effect of education and socio-economic status on social network size is reduced significantly.  This has the curious implication that traditional network theory may more accurately describe the process of cultural taste formation among low cultural capital individuals subject to high levels of network constraint \citep{burt05} while the process of culture conversion may be a more accurate representation of the deployment and operation of cultural taste among culturally advantaged elites.  Consistent with this suggestions, a recent study by Lizardo \citeyearpar{lizardo11}, shows that the extent to which individuals specialize on a given cultural pursuit is associated with the extent to which they are likely to have contacts in their ``core discussion network'' who are also acquainted with one another.
